\section{Iteration, rounding, and computation}\label{sec:7}

The two geometric threads also clarify what the result contributes to a cutting-plane method. The contraction theorem gives an exact potential decrease. The certificate gives an instance-dependent bound and a stopping criterion for an ellipsoid subproblem. Rounding supplies a controlled coordinate system for subsequent computations.

\subsection{The exact-center iteration bound}\label{sec:7.1}

\begin{corollary}[{successive exact central cuts}]\label{cor-iteration}
Suppose
$K_{j+1}=K_j\cap H_j$, where each cutting hyperplane passes through $c(K_j)$. Then
\begin{equation}\label{eq:7.1}
w(K_k)\le R_n^k w(K_0).
\end{equation}

If $K_0$ is contained in a ball of radius $R$ and all retained regions contain a target ball of radius $r>0$, then
\begin{equation}\label{eq:7.2}
k\le\frac{n\log(R/r)}{-\log R_n}.
\end{equation}

\end{corollary}

\begin{proof}
Iterate the contraction, then compare the resulting upper bound
with the positive volume that the retained target ball requires.
Induction using Theorem~\ref{thm-sharp} gives \eqref{eq:7.1}. Monotonicity of $w$ and the two ball inclusions imply
\begin{equation}\label{eq:7.3}
\kappa_n r^n\le w(K_k)\le R_n^k w(K_0)\le R_n^k\kappa_n R^n.
\end{equation}

Since $\kappa_n>0$ and $0<R_n<1$,
\begin{equation}\label{eq:7.4}
n\log r\le k\log R_n+n\log R
\quad\Longrightarrow\quad
k(-\log R_n)\le n\log(R/r),
\end{equation}
which is \eqref{eq:7.2}. 
\end{proof}

The product bound treats successive cuts separately. A natural direction is to investigate whether compatibility of near-equality geometries across recentering yields a sharper bound for consecutive cuts.

Thus an oracle process satisfying the retained-ball assumption cannot continue producing such cuts beyond this bound. For minimization, an appropriate objective-sublevel inradius must be established before turning \eqref{eq:7.2} into an objective-accuracy guarantee. This qualification is the same one that enters other volume-based localization analyses.

The exact constant improves the finite-dimensional coefficient, while preserving the known $O(n\log(R/r))$ order of the inscribed-ellipsoid method. Its dimension dependence is explicit:
\begin{equation}\label{eq:7.5}
-\log R_n=\log2-\frac12+
\sum_{j=1}^{\infty}\frac1{2j(j+1)n^j}
=\log2-\frac12+\frac1{4n}+\frac1{12n^2}+\frac1{24n^3}+O(n^{-4}).
\end{equation}

For $n\ge2$, the absolutely convergent logarithmic series gives
\begin{equation}\label{eq:7.6}
\begin{aligned}
-\log R_n
&=\log2+\frac{n-1}{2}\log(1-1/n)\\
&=\log2-\frac12
+\sum_{j\ge1}\left(\frac1{2j}-\frac1{2(j+1)}\right)n^{-j}.
\end{aligned}
\end{equation}

This proves \eqref{eq:7.5}; its positive series terms decrease with $n$, so $R_n\uparrow\sqrt e/2$. Together with $R_1<R_2$, the monotonicity holds in every positive dimension.

\subsection{Rounding also helps initialize the next subproblem}\label{sec:7.2}

The inclusion \eqref{eq:1.2} makes the value of computing a John ellipsoid broader than center selection. In its coordinates, $B_2^n\subseteq K\subseteq nB_2^n$, so both a feasible inner ball and a controlled outer radius are available.

After a central cut, the ball with center $-v/4$ and radius $1/8$ remains feasible. Translating and scaling therefore give the explicit bound
\begin{equation}\label{eq:7.7}
B_2^n\subseteq 8\bigl((K\cap H_v)+v/4\bigr)
\subseteq(8n+2)B_2^n.
\end{equation}

For $\|z\|\le1$, the point $x=-v/4+z/8$ satisfies
\begin{equation}\label{eq:7.8}
\|x\|\le\frac38<1,\qquad v^Tx\le-\frac14+\frac18=-\frac18<0,
\end{equation}
proving the first inclusion. For $x\in K\cap H_v$, the rounding bound gives
$\|8(x+v/4)\|\le8\|x\|+2\le8n+2$, proving the second.
Thus an exact John normalization supplies an explicit next-subproblem initialization with radius ratio $O(n)$. For repeated computation, natural extensions include feasible warm starts and certified cuts through approximate centers.

\subsection{Computable bounds for individual instances}\label{sec:7.3}

The universal theorem bounds the worst case. A particular instance can admit a stronger certificate. One may fix the balanced weights in \eqref{eq:3.5} and optimize only the vector multipliers, or optimize all multipliers in the normalized dual. The explicit certificate is feasible for both choices, so such optimization can only improve its determinant bound.

For any polytope ellipsoid subproblem, a primal--dual feasible pair with gap $G$ satisfies
\begin{equation}\label{eq:7.9}
\operatorname{vol}\mathcal E(c,A)\ge e^{-G}w(\mathcal P).
\end{equation}

Indeed, with $p_*:=\log(w(\mathcal P)/\kappa_n)$, weak duality gives
\begin{equation}\label{eq:7.10}
p_*\le-\log\det M=\log\det A+G
\quad\Longrightarrow\quad
\frac{\operatorname{vol}\mathcal E(c,A)}{w(\mathcal P)}
=e^{\log\det A-p_*}\ge e^{-G}.
\end{equation}

This is a relative-volume stopping criterion for a primal--dual feasible computation. Relating this gap to center error could lead to certified contraction guarantees for cuts through computed centers.

For a fixed John frame and cut direction, let $\mathfrak D(\mathcal U,v)$ denote the maximum determinant in the normalized dual. Exact duality and the contact reduction give
\begin{equation}\label{eq:7.11}
\frac{w(P_{\mathcal U}\cap H_v)}{\kappa_n}
=\mathfrak D(\mathcal U,v)^{-1},\qquad
C_n=\left(\inf_{\mathcal U,v}\mathfrak D(\mathcal U,v)\right)^{-1},
\qquad \inf_{\mathcal U,v}\mathfrak D(\mathcal U,v)=R_n^{-1}.
\end{equation}

The last infimum is attained by the equality models. Exact duality and
contact recovery were proved in Section~\ref{sec:2.3}. Section~\ref{sec:10} discusses further
properties of the certificate and related matrix inequalities.
